\section{Introduction}
\label{sec:intro}

Line heating is a production process used in shipyards to bend steel plates by moving a localized heat source along prescribed paths. The steep thermal gradients generated by the heating and cooling cycle produce permanent curvature, enabling hull, shell, and keel plates to be shaped without dedicated forming dies. Although the process is well established in practice, quantitative planning remains difficult because the final shape depends on coupled thermal diffusion, boundary cooling, material softening, plastic or inherent strain accumulation, and the order in which heating passes are applied.

For ship-production use, a simulation tool must do more than generate visually plausible temperature fields. It should accept manufacturing variables such as heating-line position, energy per unit length, travel speed, pass count, and quench schedule; predict residual deformation after cooldown; and provide outputs that can be compared with measurements or target camber profiles. This paper therefore focuses on a reproducible finite element workflow for forward line-heating prediction, with an emphasis on validation and process-scheduling effects rather than on a fully coupled electromagnetic or elastoplastic constitutive model.

\subsection{Literature Review}
Analytical heat-source descriptions remain foundational for welding and localized thermal processing. Rosenthal's moving point/line source solutions provide a baseline understanding of temperature fields and scaling trends in traveling heat inputs \cite{rosenthal1941}. For practical finite element implementation, distributed heat-source models are widely used; Goldak's double-ellipsoidal formulation is a canonical choice that enables calibration to measured melt-pool or thermal histories \cite{goldak1984}.

In ship plate forming, thermo-mechanical FEM has been used to predict temperature fields, distortion, and residual deformation under line-heating schedules. Early 3D simulations demonstrated the feasibility of predicting plate distortion under torch heating and highlighted the sensitivity to process parameters and boundary conditions \cite{biswas2007}. For large structures, computational cost motivates reduced-order strategies and decomposition methods; iterative and substructuring approaches have been proposed to accelerate prediction of welding-induced deformation while retaining key physics \cite{murakawa2015}. Validation of thermal histories and residual stresses in related multi-pass thermal processes is commonly performed via detailed transient FEM and comparison to experiments \cite{deng2012}.

Beyond forward prediction, inverse planning seeks heating parameters and paths that produce a target curvature. Data-driven surrogates and hybrid optimization have recently been explored for inverse design of ship plate forming, for example combining deep neural networks with genetic algorithms to map target geometry to process parameters \cite{yang2025}. In parallel, induction heating is an established manufacturing route where electromagnetic-thermal coupling can be modeled in detail; such literature provides context for the reduced full-line thermal-flux approximations considered here \cite{rast2018}. The present work does not claim to solve the surrounding coil field. Instead, it evaluates whether a full-line heat-input approximation can reproduce the forming response when the physical pass schedule is preserved.

\paragraph{Contributions}
\begin{itemize}
  \item A documented three-dimensional thermo-mechanical FEM workflow for ship-plate line heating, including moving-torch and full-line thermal-flux source options.
  \item Validation against an available subset of Li et al.\ (2023) Q345 benchmark cases, with calibration-anchor and non-anchor validation metrics reported separately.
  \item A process-scheduling comparison showing that collapsed multi-pass heating can substantially underpredict residual deflection, whereas sequential passes recover the benchmark response.
  \item Journal-style post-processing of FEM surfaces, residual profiles, parity plots, and process-comparison figures suitable for reproducible reporting.
  \item A keel-plate application demonstration that illustrates how the validated forward workflow can support ship-production planning.
\end{itemize}

\paragraph{Reproducibility}
All simulation inputs, post-processing scripts, and output datasets are archived in the accompanying project repository. The appendix summarizes the commands needed to reproduce the validation tables and figures.
